\chapter{General Discussion}\label{ch:discussion}

\section{Introduction}\label{sec:ch7_introduction}

This thesis asked how routine genomic surveillance could identify pairs of sampled infections compatible with recent linkage and use those pairwise relationships to investigate possible superspreading. The scenario introduced in Chapter~\ref{ch:introduction} made the difficulty clear: similar genomes sampled close together in time may suggest a shared exposure, but several transmission histories can produce that pattern. Chapter~\ref{ch:literature} established the biological and observational reasons for this ambiguity. The subsequent chapters developed a method for assessing pairwise compatibility, applied it to Scottish surveillance data, screened for unusual cluster patterns, and examined the populations associated with them.

The combined contribution is a structured way to move from genomic observations to questions for epidemiological investigation. The Scottish findings show why the stages belong together: the meaning of a candidate cluster depends on the geographic, temporal, and surveillance background in which it appears. This chapter brings those findings together around the three research questions, considers their implications, and identifies the remaining requirements for validation and practical use.

\section{Synthesis of the Main Findings}\label{sec:ch7_synthesis}

\subsection{Assessing recent linkage from genomes and sampling dates}\label{sec:ch7_synthesis_linkage}

The first research question concerned how genetic and temporal information could be combined to assess recent linkage. Chapter~\ref{ch:epilink} addressed this through EpiLink. It compared observed genetic distances and sampling-time differences with simulations of direct transmission or infection from a shared source. The simulations represented variation in infection timing, testing delay, and mutation accumulation under chosen assumptions. This gave pairwise relationships an explicit interpretation without requiring labelled transmission pairs for model fitting. The resulting compatibility weights could then be used to construct graphs and identify clusters.

The evaluation supported this use while defining its limits. In the corrected baseline synthetic benchmark, ESD was the best-performing EpiLink configuration, with BCubed F1 of 0.601 compared with 0.688 for logistic regression trained on labelled pairs. The comparator also distinguished linked pairs more accurately, and dated-tree TreeCluster achieved higher F1 than the EpiLink variants on both synthetic input types. EpiLink also recovered clusters containing high proportions of cases linked to documented Boston outbreaks. These results support using the scores to identify groups relevant to outbreak investigation, while leaving individual transmission pathways uncertain. The methodological contribution is therefore a defined and testable measure of compatibility with selected scenarios. It is not a transmission probability or a guarantee that the same performance will hold in other surveillance settings.

\subsection{Identifying unusual cluster patterns through time}\label{sec:ch7_synthesis_screening}

The second research question concerned whether the temporal patterns of recent-linkage clusters could reveal signals compatible with concentrated transmission. Chapter~\ref{ch:genomic_networks} first established the background for this question. Using EpiLink's ES formulation, which allows random variation in mutation counts, the Scottish analysis constructed compatibility graphs from 322,366 high-quality genomes within rolling windows and viral lineages. Compatibility weight was more concentrated within shared residential categories (health board, local authority, and urban/rural class) than within age, sex, or deprivation groups. Measures of graph size and connectivity were more closely associated with the numbers of sequenced genomes and positive tests than with sequencing coverage. Large or geographically concentrated clusters therefore needed to be assessed against the conditions in which they were observed.

Chapter~\ref{ch:cluster_transition_screening} used this context to screen clusters along two dimensions. The local burst score combined cluster size relative to its sampled background with the proportion of members absent from directly linked earlier clusters, where those links were observed. Here, `local' refers to the cluster and its links in the transition graph. Onward burden described additional sampled infections in later linked clusters relative to the starting cluster's size.

Among 8,962 transition-graph clusters containing at least six sequences, 631 received candidate superspreading-compatible (CSC) status: 428 for local burst only, 191 for onward burden only, and 12 for both. The small overlap shows that the two screens selected largely different clusters. It supports considering the appearance of large or novel sampled clusters separately from additional sampled infections in later linked clusters when deciding which patterns deserve investigation.

Together, Chapters~\ref{ch:genomic_networks} and~\ref{ch:cluster_transition_screening} show why a cluster's size or location alone is insufficient to establish that it is unusual. Compatibility links already tend to connect samples with shared residential attributes, and the number of observed genomes changes with both epidemic size and surveillance activity. Comparing each cluster against other eligible clusters observed under comparable conditions provides a more explicit basis for screening. The result is a set of candidates whose selection can be traced to stated rules, rather than an estimate of the number of superspreading events in Scotland.

\subsection{Understanding the populations represented in the signals}\label{sec:ch7_synthesis_characterisation}

The third research question concerned demographic, socioeconomic, and geographic patterns in the inferred graphs and clusters. Chapter~\ref{ch:candidate_characterisation} showed that CSC clusters had identifiable, though generally modest, population differences. Young adults aged 15--24, males, and urban residents made up larger shares of sampled records in CSC clusters than in background clusters. These were sequence-window records, so a sequence could contribute in more than one window. Representation did not increase or decrease consistently across deprivation groups. These comparisons describe the composition of sampled clusters, rather than individual infection risk. They do not imply that deprivation was unimportant to infection risk or to inequalities in the wider epidemic.

The distinction between composition and diversity was central to interpreting these profiles. Composition described which groups were represented; diversity described how mixed those groups were within a cluster. Large clusters could include several groups while remaining concentrated in one of them. Comparisons with random samples of the same size from the same surveillance window helped distinguish this concentration from the greater opportunity for larger clusters to contain several groups. In models using diversity relative to the simulated reference for the same cluster size and surveillance window, lower age and health-board diversity was associated with CSC status and higher local burst scores.

The geographic findings connect all three Scottish chapters. Chapter~\ref{ch:genomic_networks} found that compatibility weight was more concentrated within shared residential categories than within age, sex, or deprivation groups. Chapter~\ref{ch:cluster_transition_screening} distinguished local bursts from onward burden. In Chapter~\ref{ch:candidate_characterisation}, lower health-board diversity relative to the simulated reference was associated with higher local burst scores, whereas greater diversity on the same scale was associated with higher onward burden scores. Taken together, these results motivate separate investigations of concentrated local exposure and connections between regions. They do not establish those mechanisms. The local burst score includes cluster size, whereas onward burden is measured relative to that size and is scored only for clusters with additional sampled infections in later linked clusters. In addition, links in the transition graph record shared sampled membership rather than transmission between clusters.

The population findings also remained tied to epidemic and surveillance context. The models estimated differences between policy periods and clades, and adding measured incidence and sequencing variables changed most CSC associations little. That stability supports further investigation of the observed profiles, but does not remove the effects of unmeasured differences in which infections were tested and sequenced. These results provide evidence of differences in the population profiles of sampled clusters. Establishing the causes of those differences requires information beyond the surveillance record analysed here.

\section{Implications for Genomic Surveillance}\label{sec:ch7_triage_layer}

The findings suggest a role for genomic screening in prioritising epidemiological review. A candidate could be presented with its local burst and onward burden scores, sampling dates, lineage, population composition, and surveillance context. Investigators could then compare that information with contact-tracing and outbreak records. For example, a burst candidate concentrated in a particular residential area or population group could prompt examination of a shared activity or institution, while a candidate spanning several residential areas with a high onward burden score could prompt examination of links between communities. These are proposed uses of the findings; the thesis did not test whether such review improves public health decisions.

Examining population attributes after assigning candidate status made it possible to describe the signals without selecting them directly on those attributes. It also preserved a useful distinction between identifying an unusual genomic pattern and explaining why it occurred. Population and place nevertheless influence which infections enter surveillance and how genomic relationships are distributed. Interpretation therefore needs to retain both the candidate's observed profile and the background against which it was selected. The practical value of the approach will depend on whether these combined descriptions yield useful, timely leads at an acceptable review workload.

\section{Strengths and Limitations}\label{sec:ch7_limitations}

A strength of the thesis is the connection between methodological evaluation and a national empirical application. EpiLink supplied a defined pairwise measure; the Scottish analyses examined its graph structure, screened clusters over time, and characterised the resulting signals. The linked data allowed those signals to be considered across changing incidence, policy, testing, vaccination, and viral lineages. This made surveillance context part of the interpretation throughout the analysis.

The main limitation is incomplete observation. The data contain only infections that were detected, successfully sequenced, and retained after quality control and linkage. Calibration and regression adjustment compare observed records; they cannot recover missing infections or fully correct for targeted outbreak sequencing. Real outbreaks may be missed, split across clusters, or merged with other sampled infections. Conversely, an unusual cluster may reflect intensive sampling or repeated introductions. Without an independent reference set of Scottish superspreading events, the sensitivity, specificity, and positive predictive value of CSC status remain unknown. Passing the screening threshold means that a score is unusual under the reference distribution used to calibrate the screen. It does not give the probability that a candidate corresponds to a true superspreading event.

Uncertainty also arises from the model and graph construction. The synthetic benchmark used a single fixed transmission tree, with the same process-based model used to generate and analyse the synthetic data. Performance fell when genetic resolution was reduced and under some parameter mismatches. The later analyses depended on the selected linkage scenarios, temporal windows, lineage groups, the threshold for removing weak edges, the Leiden resolution parameter, and screening rules. Links in the transition graph represented shared sampled membership, so sampling gaps, window boundaries, and changes in cluster assignments could alter onward burden. The final population estimates did not fully account for uncertainty introduced by the earlier analysis stages.

Characterisation introduced further limits. Recorded residence did not establish where infection occurred. Policy periods and clades reflected related changes in the epidemic. The simulated reference accounted for cluster size and the sampled population's composition within each window, but did not match clusters by lineage or preserve associations between attributes. In the composition models, larger clusters contributed more records, members shared a cluster outcome, and sequences could recur across windows. This dependence was not fully modelled, so uncertainty intervals may be too narrow. Both local burst models and the expanded onward burden model fitted to sequence-window records retained warnings in the Bayesian sampling diagnostics. Checks of how well these models reproduced the observed score distributions were not reported. These issues favour interpreting broad population patterns over treating individual estimates as definitive.

\section{Priorities for Further Work}\label{sec:ch7_future_work}

The first priority is external epidemiological validation. Burst-only, burden-only, and joint candidates should be compared with independently assembled outbreak and contact-tracing records. Assessing missed outbreaks would also require examining recorded events that were not selected by the screen. Validation should allow an outbreak to overlap several genomic clusters and a cluster to include more than one recorded event. This would help distinguish useful epidemiological signals from patterns created mainly by sampling or analytical choices.

The second priority is to test how changes at each analysis stage affect candidate selection and population findings. This includes changes in biological timing parameters, the inclusion of unsampled intermediates, window design, thresholds for removing weak edges, and clustering rules. Retaining most compatibility weight when weak edges are removed is not sufficient evidence that the same clusters or candidates are recovered. Further work should account for dependence within composition models, address the Bayesian sampling warnings, and report checks of how well the fitted models reproduce the observed score distributions. The aim is to identify which findings persist across plausible alternative modelling and screening choices.

The third priority is to evaluate the screen in practice as new data become available. This should measure the time taken to generate signals, repeated alerts concerning the same investigation, reviewer workload, and the additional information obtained beyond existing surveillance. Such evaluation would assess whether the proposed screening process helps investigators make better decisions and which computational or reporting changes are needed before routine use.

\section{Conclusion}\label{sec:ch7_conclusion}

This thesis developed a method for assessing recent-linkage compatibility and connected it to the screening and characterisation of national genomic data. The Scottish analyses showed that geographic structure and surveillance conditions are essential to interpreting unusual clusters. They also showed that the local burst and onward burden screens selected largely different candidates, while their scores had contrasting associations with health-board diversity relative to the simulated reference. These findings provide specific questions about local concentration and connections between populations that can be tested with epidemiological records.

The overall contribution is a reproducible route from genomic relationships to focused investigation, with a clear account of what each stage measures. Its practical value now depends on external validation and evidence that the resulting signals improve surveillance decisions.
